\section{Introduction}
\label{sec:introduction}

The $k$-core of a graph is its largest induced subgraph in which every vertex
has degree at least $k$ within that subgraph. This work focuses on the
$3$-core because, in classical random graph models, increasing connectivity
can cause it to appear abruptly. In the large-graph limit, the proportion of
all vertices belonging to the core, called the core fraction, can jump from
zero to a positive value~\cite{janson2007kcore}. Near this transition,
finite graphs of the same size and generated at the same parameter value may
contain almost no core or retain a substantial fraction of their vertices.

This work studies how this fraction varies between finite random
intersection graphs (RIGs) built from overlapping triangle communities.
The aim is to understand and predict both how often a substantial $3$-core
appears near the transition
and how its fraction varies among the graphs that retain one. Here, survival
means that the core fraction reaches a fixed positive cutoff, which
distinguishes substantial cores from small nonempty ones.

In the Poisson triangle RIG studied in this work, each individual receives
a Poisson number of community memberships, which are randomly grouped in threes.
Individuals sharing a community are connected, so these groups usually form
triangles, with overlaps where individuals belong to several communities
\cite{hofstad2021communities}. Thus, the Poisson parameter controls average
connectivity in this model. The edges between neighbours create local clustering
and allow the neighbours to support each other during core pruning. Whether one
remains in the $3$-core can therefore depend on whether the other remains
\cite{gleeson2009clustered,bhat2017clustered}.

For comparison, a second graph model, the configuration model (CM), is also
studied. Its degrees are chosen to match the RIG's limiting degree
distribution, and edges are formed by randomly pairing the corresponding
edge ends. Connections are therefore formed without organizing vertices
into triangle communities. Although degrees in individual finite graphs
need not coincide, the common limiting distribution allows the comparison
to examine how arranging connections into communities affects core emergence
and size beyond what
the limiting degree distribution explains. These effects include several
dependencies introduced by communities and cannot be attributed to
clustering alone. Established theory also describes the CM's core fluctuations
near the transition~\cite{janson2008fluctuations}. This provides a starting
point for investigating whether a similar description applies to the
triangle RIG.

Motivated by these results, this work investigates an approximation for
how often a substantial $3$-core appears in finite graphs near the
transition and how its fraction is distributed among graphs that retain
one. The approximation represents differences between graphs by treating
each graph as if it had a slightly different connectivity threshold for
core survival~\cite{janson2008fluctuations}. A substantial core is predicted
when connectivity exceeds this threshold, and its predicted fraction
depends on how far above the threshold the connectivity lies. The
approximation is constructed separately for the RIG and CM, using each
model's theoretical relation between connectivity and core fraction.
The numerical study examines how the survival probability and the
distribution of core fractions among surviving graphs change with graph
size. The central aim of this numerical study is testing
whether theory and observations of how often
cores survive are sufficient to predict that distribution in an
independent set of simulations. Comparing the approximation's accuracy
for the two models assesses how well the description supported by CM
theory extends to the triangle RIG.

\Cref{sec:models-and-observables} defines the two graph models and the
measures of core survival and size. \Cref{sec:theoretical-predictions}
develops the theoretical predictions for the transitions and finite-size
fluctuations. \Cref{sec:numerical-methodology} explains the simulation
design and how the predictions are evaluated, followed by the results and
model comparisons in \cref{sec:numerical-results}. \Cref{sec:discussion}
interprets these results and discusses the limitations of the
approximation, while \cref{sec:conclusion} summarizes the main conclusions.
Appendices~\ref{app:derivations-and-checks} and~\ref{app:experiment-details}
provide supporting derivations, numerical checks, and details of the
experiments.
